% !TEX root = ../../../main.tex

\subsection{模型发布、重载与实时识别联动}


\WhuAutoFigure{0.82\textwidth}{0.42\textheight}{figures/chapter06/fig6-4-model-release-reload-flow.png}{模型发布、重载与实时识别联动流程图}{fig:chapter06-model-release-reload-flow}


固定词表模型训练完成后，训练产物还需要进入系统运行链路，才能真正被移动端实时识别功能使用。
HearBridge 系统通过模型版本管理、模型发布和识别服务重载机制，将训练产物接入后台管理端和实时识别服务。

模型发布过程中，管理员在后台管理端选择某个模型版本并点击发布。
管理端向 Spring Boot 服务端发送发布请求，服务端根据模型版本 ID 查询对应记录，并检查模型文件和标签映射文件是否可用。
随后，服务端调用 Python 手势识别服务提供的重载接口，使 Python 服务加载指定模型文件和标签映射。

实时识别过程中，Python 服务根据当前加载的模型输出预测类别和置信度，并将结果通过 WebSocket 返回移动端。
返回结果中除了识别标签和置信度外，还可以包含当前模型版本名称、是否使用发布模型等状态信息。
